\section{A point-only synthetic geometry}
\label{sec:synthetic}
Let $K$ denote the effective ordered real-closed field of real algebraic
numbers. In this section $A,B,C,\ldots$ are elements of $K^2$. Write
$u\cdot v=u_1v_1+u_2v_2$, $N(u)=u\cdot u$,
$[u,v]=u_1v_2-u_2v_1$, and $Ju=(-u_2,u_1)$. Then
$N(u)=0$ precisely when $u=0$, since squares are nonnegative and a
nonzero coordinate has strictly positive square. Put
$|u|=\sqrt{N(u)}$, using the nonnegative root.
We verify the usual plane axioms in the point language of Tarski.
The forms of the ten individual axioms below agree with the GeoCoq
presentation of the Schwabh\"auser--Szmielew--Tarski system.%
\footnote{GeoCoq, \emph{Tarski axioms},
\url{https://geocoq.github.io/GeoCoq/html/GeoCoq.Axioms.tarski_axioms.html}.}

\paragraph{Two geometric relations.}
Write $AB\equiv CD$ when $N(B-A)=N(D-C)$. Write $B(A,X,C)$ when
$X=(1-t)A+tC$ for some $0\leq t\leq1$. There is a quantifier-free test
and there is a computable choice of $t$. With $u=C-A$ and $v=X-A$,
the test is
\[
 [u,v]=0,\qquad v\cdot(v-u)\leq0 .
\]
When $u=0$ the inequality forces $v=0$ and we take $t=0$.
Otherwise choose the first nonzero coordinate $u_i$ and put $t=v_i/u_i$.
The determinant equation gives $v=tu$, and the inequality becomes
$t(t-1)N(u)\leq0$, precisely $0\leq t\leq1$. The converse follows
by substitution, including when endpoints coincide. Equality and both
relations are decidable on point codes. Primitive recursive realizer
formulas below branch only on decidable comparisons.

\begin{theorem}[The ten point axioms]
The exact plane satisfies segment-congruence reversal, common-segment
transitivity, identity of congruence, segment construction,
the five-segment axiom, identity of betweenness, inner Pasch,
both plane-dimension axioms, and Tarski's Euclidean axiom.
\end{theorem}

\begin{proof}
\emph{Congruence and betweenness identities.}
The equality $N(-u)=N(u)$ gives $AB\equiv BA$.
From $AB\equiv CD$ and $AB\equiv EF$ ordinary transitivity gives
$CD\equiv EF$. From $AB\equiv CC$ one gets $N(B-A)=0$, hence $A=B$.
Any affine combination of $A$ with itself is $A$, so $B(A,X,A)$ gives $X=A$.

\emph{Segment construction.}
Given $A,B,C,D$, put $d=|D-C|$. The goal is $B(A,B,E)$ and $BE\equiv CD$.
If $A=B$, take $E=B+(d,0)$ and use parameter zero for betweenness.
Otherwise $l=|B-A|>0$. Set $E=B+(d/l)(B-A)$. Now
$N(E-B)=d^2$ and $B=A+\tfrac{l}{l+d}(E-A)$ with
$0<l/(l+d)\leq1$. In particular, a requested segment of length zero
is allowed.

\emph{Five segments.}
Suppose $B(A,B,C)$, $B(A',B',C')$, $A\ne B$, and the four pairs of
corresponding segments $AB$, $BC$, $AD$, $BD$ are congruent.
Let $u=B-A$, $w=D-B$. Writing $B=A+t(C-A)$ gives $t>0$, so
$C=B+ku$ where $k=(1-t)/t\geq0$. Define primed quantities likewise,
using $N(u')=N(u)>0$. Since $k^2N(u)=N(C-B)=k'^2N(u')$,
nonnegativity gives $k=k'$. Moreover,
\[
  2u\cdot w=N(D-A)-N(u)-N(w).
\]
All entries on the right agree with their primed counterparts. Thus
$N(D-C)=N(w)-2k(u\cdot w)+k^2N(u)$ agrees with $N(D'-C')$, as required.
This proof permits $B=C$.

\emph{Inner Pasch.}
Given $B(A,P,C)$ and $B(B,Q,C)$, write
$P=(1-p)A+pC$, $Q=(1-q)B+qC$, with $p,q\in[0,1]$.
Set $h=p+q-pq=p+(1-p)q=q+(1-q)p$. If $h=0$,
then $p=q=0$, so $P=A$, $Q=B$ and $X=A$ works. If $h>0$,
put $r=p(1-q)/h$ and $s=q(1-p)/h$. The identities
$1-r=q/h$ and $1-s=p/h$ show $r,s\in[0,1]$. Expansion gives
\[
  (1-r)P+rB
  =\frac{q(1-p)A+p(1-q)B+pqC}{h}
  =(1-s)Q+sA .
\]
Use this point for $X$. Thus $B(P,X,B)$ and $B(Q,X,A)$ hold even
when the given triangle is collinear.

\emph{Dimension.}
The points $(0,0),(1,0),(0,1)$ have determinant $1$, so none of the
three cyclic betweennesses holds. Conversely let $P\ne Q$ and suppose
each of $A,B,C$ is equidistant from $P$ and $Q$. The locus satisfies
\[
                2X\cdot(Q-P)=N(Q)-N(P).
\]
Writing $u=Q-P$, the differences $B-A,C-A$ lie in the span of $Ju$.
Indeed $u,Ju$ form an orthogonal basis of determinant $N(u)>0$.
Thus $A,B,C$ have respective line parameters $0,b,c$.
Decidably sort these three values. The point with median parameter
lies between the other two: for $a\leq m\leq z$ use $(m-a)/(z-a)$
when $a<z$, and use zero when all three agree. Reversing endpoints
preserves betweenness by replacing $t$ by $1-t$. Consequently
$B(A,B,C)\lor B(B,C,A)\lor B(C,A,B)$.

\emph{Euclidean axiom.}
Assume $B(A,D,T)$, $B(B,D,C)$ and $A\ne D$. Write
$D=A+s(T-A)=(1-t)B+tC$ with $s,t\in[0,1]$. Since $D\ne A$,
$s>0$. Set
\[
 X=A+(B-A)/s,\qquad Y=A+(C-A)/s .
\]
Then $B=A+s(X-A)$, $C=A+s(Y-A)$ and
$T=A+(D-A)/s=(1-t)X+tY$. These are precisely
$B(A,B,X)$, $B(A,C,Y)$ and $B(X,T,Y)$.
\end{proof}

\begin{theorem}[Elementary continuity, with point parameters]
Let $L(x)$ and $U(y)$ be any fixed first-order point formulas, allowing
parameters. If there is an $A$ such that $B(A,x,y)$ for every
$x\in L,y\in U$, then there is an $S$ such that $B(x,S,y)$ for every
$x\in L,y\in U$. For each instance $S$ is obtainable by primitive
recursive operations on the parameter codes and the given $A$.
\end{theorem}
\begin{proof}
Translate point variables into pairs of field variables, using the
polynomial tests above. The effective real-closed-field decision and
sampling procedure decides emptiness of each definable set and produces
a member in the inhabited case. If either $L$ or $U$ is empty, take $S=A$.
Decide whether some $P\in L$ differs from $A$. In the negative case
every $x\in L$ equals $A$; again $S=A$ works. In the positive case
choose such a $P$ and some $Q\in U$. Put $u=P-A$. From $B(A,P,y)$
each $y\in U$ has the form $A+bu$ with $b\geq1$. From $B(A,x,Q)$
each $x\in L$ has the form $A+au$ with $a\geq0$. The hypothesis
$B(A,x,y)$ gives $a\leq b$ for every such pair. Thus
\[
             I=\{a\in K:L(A+au)\}
\]
is inhabited and has an upper bound. After quantifier elimination
the signs of finitely many univariate polynomials decide membership.
Discard identically zero and constant polynomials after recording
their signs. The ordered list of all distinct real roots partitions
the line into singleton and open-interval cells, with constant truth
on each cell. Choose the last occupied cell. It cannot be right-unbounded.
Its unique point or its right endpoint is $c\in K$.
Every member of $I$ is $\leq c$, and $c$ is at most every upper bound:
in the open-cell case, if an upper bound were $b<c$, the interval
contains a point larger than $b$ (take the midpoint of $c$ and the
larger of $b$ and any sample in that interval). Hence $a\leq c\leq b$
for every original pair. The one-dimensional parameter calculation in
the dimension proof gives $B(x,A+cu,y)$. This constructs $S=A+cu$.
The emptiness, sampling, sign tests, root operations and final
finite scans are those of the effective field presentation. The
schema ranges over the point language; arbitrary arithmetic predicates
on coordinate codes need not be invariant under this translation.
\end{proof}

\subsection{Lengths, angle comparison and rigid congruence}
For nonzero vectors $u,v$ the pair
\[
 (c,s)=\left(\frac{u\cdot v}{|u|\,|v|},
                  \frac{[u,v]}{|u|\,|v|}\right)
\]
encodes the directed angle from $u$ to $v$. Indeed the polynomial identity
$(u\cdot v)^2+[u,v]^2=N(u)N(v)$ gives $c^2+s^2=1$, and
$R(c,s)(u/|u|)=v/|v|$, where
$R(c,s)(x,y)=(cx-sy,sx+cy)$. These equations follow by expanding
the two coordinates and determine $c,s$ uniquely. Multiplication of
the two-by-two rotation matrices gives the angle addition law.
For ordinary unsigned angles, take $(c,|s|)$ and order from zero to
a straight angle by decreasing $c$. On the upper semicircle $c$
uniquely determines $|s|=\sqrt{1-c^2}$, so equality and order of
angles reduce to scalar comparisons. No transcendental angle unit is needed.
For example, for a noncollinear triangle let $u=B-A$, $v=C-A$.
The square of the length opposite the angle at $A$ is
\[
         N(C-B)=N(u)+N(v)-2|u|\,|v|c .
\]
Thus the Pythagorean relation is equivalent to perpendicularity
$u\cdot v=0$. The formula also gives side-angle-side from the
three-side congruence theorem below.

The same polynomial identity yields $u\cdot v\leq |u|\,|v|$.
Expansion of $N(u+v)$ and comparison of nonnegative square roots now
give $|u+v|\leq |u|+|v|$. Equality requires
$u\cdot v=|u|\,|v|$, equivalently $[u,v]=0$ and $u\cdot v\geq0$.
If $u,v$ are both nonzero, these conditions say $v=tu$ for $t\geq0$:
choose a nonzero coordinate of $u$ and use the determinant equation.
Applied to $u=B-A$, $v=C-B$, the equality case is precisely $B(A,B,C)$.
When either vector is zero the endpoint cases give the same conclusion.

\begin{proposition}[From three congruent sides to a motion]
Suppose the three corresponding sides of the ordered triples $A,B,C$ and
$A',B',C'$ are congruent. A translation followed or preceded by an
orthogonal linear map carries the first triple to the second. A code
for such a map and translation can be computed from the six points.
\end{proposition}
\begin{proof}
If $A\ne B$, let $l=|B-A|$, $e=(B-A)/l$ and $f=Je$;
define $e',f'$ using the equal primed length. These are orthonormal
frames, and $C-A=xe+yf$, where $x=(C-A)\cdot e$, $y=(C-A)\cdot f$.
Expansion gives
\[
              2lx=l^2+N(C-A)-N(C-B).
\]
Thus $x'=x$, and $y'^2=N(C'-A')-x'^2=N(C-A)-x^2=y^2$.
If $y'=y$ put $\epsilon=1$; otherwise $y'=-y$ by
$(y'-y)(y'+y)=0$, and put $\epsilon=-1$. The map
$A+he+kf\mapsto A'+he'+\epsilon kf'$ has the claimed properties;
its inverse has the same formula with frames interchanged.
If $A=B$ but $C\ne A$, permute the triples and use $AC$ as base.
If all three points coincide, the primed triple also coincides, and
translation by $A'-A$ suffices. Orthogonality preserves squared norm
and an affine map preserves affine combinations, so these motions
preserve both geometric predicates. Conversely a motion preserves all
three sides.
\end{proof}
